\clearpage
\fancyhead[L]{\small SCHULZE DARSTELLUNG DES WELTSYSTEMS.}
\SourceLeaf{538}
\begin{center}\small Göttingische gelehrte Anzeigen. Stück 21. Seite 305..306. 1811 Februar 23.\\[3pt]\rule{.72\linewidth}{.3pt}\end{center}

\emph{Darstellung des Weltsystems. Ein Leitfaden für den Unterricht in der Astronomie auf Schulen. Abgefasst und zur Erleichterung des eigenen weitern Studiums der Sternwissenschaft, mit den nöthigsten literärischen Anmerkungen und Nachweisungen versehen von} G. L. \textsc{Schulze}, Prediger in Polenz und Ammelshayn bei Leipzig. Mit 4 Kupfertafeln. 1811. 390 Seiten in Octav. Leipzig, Baumgärtner.

Wenn gleich die tieferen Kenntnisse der Astronomie ihrer Natur nach nur das Eigenthum weniger Eingeweihten sein können, so enthält doch diese Wissenschaft einen solchen Reichthum jedem gebildeten Menschen zugänglicher und höchst interessanter Wahrheiten, dass jeder darin für das ganze Leben eine Quelle mannigfaltiger Genüsse und erhebender Ansichten finden kann, und es stets dankbar erkennen muss, wenn er früh Gelegenheit gehabt hat, sich die Wege dazu bahnen zu lassen. Ein zweckmässiger Unterricht darin für das jugendliche Alter ist daher in dieser Hinsicht etwas sehr Schätzbares, und es ist nicht zu bezweifeln, dass durch eine allgemeinere Verbreitung desselben, als bisher bei uns üblich gewesen ist, selbst manches Talent geweckt und zur künftigen weitern Ausbildung angereizt werden könne. Das vorliegende Lehrbuch hat, unsers Erachtens, eine zu dieser Absicht sehr zweckmässige Einrichtung. Es hält eine glückliche Mittelstrasse zwischen den beiden Klippen, woran man sonst so leicht scheitert, der zu abstracten abschreckenden Trockenheit und der seichten Oberflächlichkeit, die nur das Halbwissen befördert und den Aufflug des eigenen Nachdenkens lähmt. Ein besonderes Lob verdient die Nachweisung der vorzüglichsten astronomischen Schriften bei allen abgehandelten einzelnen Materien, wo Jeder, der weitere Belehrung sucht, seine Wissbegierde befriedigen kann. Bei einem Werke dieser Art mehr ins Einzelne zu gehen, verstattet der Zweck unserer Blätter nicht: aber gestehen müssen wir, dass es uns eine erfreuliche Erscheinung gewesen ist, bei einem Manne, dessen Berufsgeschäfte von ganz anderer Art sind, so gründliche Kenntnisse der Astronomie und eine so vertraute Bekanntschaft mit den besten astronomischen Schriften anzutreffen.

\clearpage
\fancyhead[L]{\small SEYFFER SUPER LONGITUDINE GEOGRAPHICA.}
\hypertarget{native-538Seyffer}{}\label{native-538Seyffer}
\begin{center}\small Göttingische gelehrte Anzeigen. Stück 78. Seite 773..775. 1811 Mai 18.\\[3pt]\rule{.72\linewidth}{.3pt}\end{center}

\emph{Super longitudine geographica speculae astronomicae regiae, quae Monachii est, ex triginta septem defectionibus Solis observatis et ad calculos revocatis nunc primum definita a} \textsc{Carolo Felice Seyffer}. 1810. 104 Seiten in Quart. München.

{\rightskip=0pt plus 1em Nicht aus 37 Sonnenfinsternissen, wie man aus dem Titel schliessen sollte, sondern aus der zu München gemachten Beobachtung Einer Sonnenfinsterniss (vom 16. Juni 1806), verglichen mit 36 Beobachtungen an andern Orten, leitet Hr. \textsc{Seyffer} in dieser Abhandlung, die aus den Denkschriften der königl. Baierischen Academie der Wissenschaften besonders abgedruckt ist, die Länge von München ab. Die verglichenen Orte sind: \emph{Rom} (zwei Beobachtungen), \emph{Padua, Mailand, Madrid, Aranjues, Pampelona, Rinderhook, Fort Orange} (beide in den vereinigten Staaten von Nordamerika), \emph{Amsterdam, Utrecht,} \SourceLeaf{539}\emph{Zürich, Ochsenhausen, Leipzig, Breslau, Ofen, Cracau, Erlau} (in Ungarn), \emph{Schweidnitz, Hamburg, Luck} (in Polen), \emph{Bourg en Bresse}, Insel \emph{Leon} (bei Cadiz), \emph{Montauban, Toulouse, Paris, Prag, Lilienthal, Reikevich} (auf Island), \emph{Göttingen, Neapel, Brünn. Berlin, Regensburg, Cremsmünster.} In so fern es nur die Längenbestimmung von München galt, wäre es freilich nicht nöthig gewesen, so viele Orte in Rechnung zu nehmen: als ein Beitrag zur Bestimmung der Längen dieser Orte selbst ist indess diese Arbeit, von welcher Hr. \textsc{Seyffer} das ganze Detail aller Rechnungen hat abdrucken lassen, schätzbar.\par}

Nachdem Hr. \textsc{Seyffer} aus denjenigen Beobachtungen, die vollständig waren, die Correctionen der Mondsbreite, der Horizontal-Parallaxe und der Summe der Halbmesser discutirt hat, findet er die Länge von München im Mittel aus 15 Vergleichungen $37^{\mathrm m}5^{\mathrm s}6$ östlich von Paris. In dem neuesten Bande der Connoissance des tems, wo München unter den Oertern mit aufgeführt wird, deren Längen Hr. \textsc{Burckhardt} von neuem discutirt hat, ist sie $37^{\mathrm m}0^{\mathrm s}$ angesetzt.

\clearpage
\fancyhead[L]{\small NOUET OBSERVATIONS EN EGYPTE.}
\hypertarget{native-539Nouet}{}\label{native-539Nouet}
\begin{center}\small Göttingische gelehrte Anzeigen. Stück 79. Seite 780..781. 1811 Mai 18.\\[3pt]\rule{.72\linewidth}{.3pt}\end{center}

\emph{Description de l'Egypte. Tome premier à Paris de l'imprimerie impériale} 1809. \raisebox{-0.3ex}{,}\emph{Mémoire premier. Observations astronomiques faites en Egypte pendant les années} VI. VII \emph{et} VIII (1798. 1799. 1800) \emph{par} M. \textsc{Nouet} \emph{astronome de la commission des sciences et arts d'Egypte.}

Die Instrumente, deren sich \textsc{Nouet} bediente, waren ein Multiplicationskreis von 25 Centimeter Durchmesser mit Decimaleintheilung (weiter hin wird auch ein Quadrant von 35 Centimeter Halbmesser erwähnt), eine Seeuhr von \textsc{Louis Berthoud} und ein achromatisches Fernrohr von \textsc{Dollond}, mit 63 Millimeter Oeffnung. In Kairo erhielt er von \textsc{Beauchamp}, welchen er daselbst traf, noch einen zweiten Chronometer. In das Detail der Beobachtungen, deren Ausbeute die Bestimmung der Breiten und Längen von 36 verschiedenen Punkten ist, weiter einzugehen, erlaubt der Raum und die Absicht unserer Blätter nicht; wir begnügen uns, nur Einiges davon auszuheben. Die südlichsten Punkte sind die Insel \emph{Philä} oberhalb der Catarracten des Nils, in $24^\circ\,1'\,34''$ Breite, und \emph{Syene} in $24^\circ\,5'\,23''$ Breite. Von letzterem Orte kennt man die Sage, dass zur Zeit der Sommer-Sonnenwende ein Brunnen bis auf den Grund erleuchtet worden sei. \textsc{Nouet} berechnet nach \textsc{Laplace}'s Formel den Zeitpunkt, wo die Schiefe der Ekliptik jener Breite gleich gewesen ist, auf 3430 Jahre vor unserer Zeitrechnung, und glaubt deshalb jener Sage ein so hohes Alter beilegen zu müssen. Allein dieser Schluss scheint doch nicht begründet genug, da ein solches Phänomen noch Statt finden konnte, wenn auch der \emph{Mittelpunkt} der Sonne dem Zenith in geringer Entfernung südlich vorbeiging. Die Längenbestimmungen sind grösstentheils chronometrisch; doch hat \textsc{Nouet} auch Gelegenheit gehabt, vier Ocultationen zu beobachten, nemlich die Bedeckung der Venus 1798 December 13 zu Ssalehhiyeh, die Bedeckung von $\delta$ im Scorpion 1799 April 21 zu Kairo, die Bedeckung der Venus 1799 November 23 ebendaselbst, und die Bedeckung von $\alpha$ im Scorpion 1800 Juli 28 zu Alexandrien; überdies noch zu Kairo eine beträchtliche Anzahl Verfinsterungen von Jupiterstrabanten. Die Abweichung der Magnetnadel zu Alexandrien im Mittel aus 26 Beobachtungen, deren Datum aber nicht beigefügt ist. fand sich $13^\circ\,6'$ N.\,W., und die Neigung im Mittel aus 12 Beobachtungen $47^\circ\,30'$ Nordl.
